\chapter{まとめ} \label{chap:summary}

本論文では，これよりファラデーの研究の全般を窺（うかが）った．
それにはチンダルの書いたものがあるから，そのままを紹介する．
ただファラデーの死後間もなく（今から言えば五十年前に）書かれたものだから，多少不完全な点もあるようである．

「アルプス山の絶頂に登りて，諸山岳の重畳するを見渡せば，山はおのずから幾多の群をなし，各々の群にはそれぞれ優れた山峯あって，やや低き諸峰に囲まるるを見る．
非常なる高さに聳ゆるの力あるものは，必ずや他の崇高なるものを伴う．
ファラデーの発見もまたそのごとく，優秀なる発見は孤立せずして，数多の発見の群集せる中の最高点を形成せるを見る．

「すなわちファラデーの電磁気感応の大発見を囲んで群集せる発見には，電流の自己感応，反磁性体の方向性，磁気指力線とその性質並びに配布，感応電流による磁場の測定，磁場にて誘導さるる現象．
これらは題目において千差万別なるも，いずれも電磁気感応の範囲に属するものである．
ファラデーの発見の第二の群は，電流の化学作用に関するもので，その最高点を占むるは，電気分解の定律なるべく，これを囲んで群集せるものは，電気化学的伝導，静電気並びに電池による電気分解，金属の接触による起電力説と電池の原因に関する説の欠点，および電池の化学説である．
第三群は光の磁気に対する影響で，こは正にアルプス山脈のワイスホルン峰に比すべきものか．
すなわち高く美しくしかも孤峰として聳ゆるもの．
第四群に属するのは反磁性の発見で，総ての物質の磁性を有することを中心として，焔（ほのお）並びにガス体の磁性，磁性と結晶体との関係，空中磁気（ただし一日並びに一年間の空中磁気の変化に関する説は完全とは言い難かるべきも）である．
［＃「である．」は底本では「である」］

「以上に列記したるは，ファラデーの発見中の最も著名なるもののみである．
たといこれらの発見なしとするも，ファラデーの名声は後世に伝うるに足るべく，すなわちガス体の液化，摩擦電気，電気鰻の起す電気，水力による発電機，電磁気廻転，復氷，種々の化学上の発見，例えばベンジンの発見等がある．
「かつまたそれらの発見以外にも，些細の研究は数多く，なお講演者として非常に巧妙であったことも特筆するに足るだろう．「これを綜合して考うれば，ファラデーは世界の生んだ最大の実験科学者なるべく，なお歳月の進むに従って，ファラデーの名声は減ずることなく，ますます高くなるばかりであろう．」と．

ファラデーの名声がますます高くなるだろうと書いたチンダルの先見は的中した．
しかし，チンダルはファラデーの最大の発見を一つ見落しておった．
それは実験上での発見ではなくて，ファラデーが学説として提出したもので，時代より非常に進歩しておったものである．もしチンダルにその学説の価値が充分に理解できたならば，チンダル自身がさらに大科学者として大名を残したに違いない．
それは何か．
前にフィリップスに与えた手紙のところで述べた電気振動が光であるという説である．
マックスウェルの数理と，ヘルツの実験とによりて完成され，一方では理論物理学上の最も基本的な電磁気学に発展し，他方では無線電信，無線電話等，工業界の大発明へと導いた光の電磁気説がそれである．
